\chapter{文献导读}
\label{app:g}

这份导读按阅读目的组织文献。正文各章保留更密集的引用；这里列出进入主题时适合先读的一组。所有条目使用 NASA ADS bib key，并在 \texttt{references.bib} 中维护。

\section{经典起点}
\label{appsec:g-classic}

强度干涉的历史起点是 Hanbury Brown 和 Twiss 的实验与天文实现 \citep{1956Natur.178.1046H,1957RSPSA.242..300B,1974iiia.book.....H}。读这组文献时，重点在于看他们怎样把相关峰、望远镜面积、基线、带宽和积分时间连成一个可执行观测，不要停在 \(g^{(2)}\) 的名字上。量子光学基础可从 Glauber 和 Sudarshan 的相干性理论，以及 Mandel--Wolf 和 Gerry--Knight 的教材进入 \citep{1963PhRv..130.2529G,1963PhRvL..10..277S,1995ocqo.book.....M,1985stop.book.....G}。

\section{强度干涉的现代复兴}
\label{appsec:g-sii}

现代复兴的主线是把 Cherenkov 阵列、大面积望远镜、数字相关器和校准星库合在一起。早期方案和仪器路径可读 \citet{2009A&A...508..531N,2017MNRAS.472.4126G}；真实天文结果从 VERITAS、MAGIC 和后续样本开始进入系统测量阶段 \citep{2020NatAs...4.1164A,2021MNRAS.506.1585Z,2022MNRAS.512.1722K,2024MNRAS.529.4387A,2024ApJ...966...28A,2025ApJ...995..191A}。读这些文章时，优先看数据率、零基线标定、背景扣除、系统误差地板和目标选择，角直径表要放回这些条件里理解。

\section{探测器、计时和事件表}
\label{appsec:g-instrument}

事件表天文学需要把光学、电子学和时间标准读在一起。AQuEye、IquEye 和光子计数仪器论文适合用来理解 ps--ns 级时间标签、绝对同步和高光子率数据流 \citep{2015SPIE.9504E..0CZ,2016SPIE.9907E..0NZ,2013NIMPA.725..187J,2020RScI...91a3108W}。读这组文献时，把探测器效率、死时间、afterpulsing、TDC 分辨率、GPS/White Rabbit 同步和数据格式单独列成表；这些量后来都会进入第 \ref{chap:e13}--\ref{chap:e14} 章的数据模型。

\section{量子估计和亚分辨测量}
\label{appsec:g-estimation}

Rayleigh curse 和 SPADE 的核心文献说明，传统分辨率判据并非所有参数估计问题的信息极限 \citep{2015Optic...2..646T,2016PhRvX...6c1033T,2017PhRvL.118g0801T,2017NJPh...19b3054T}。读这组文献时，把每篇文章的源模型、PSF、背景、可用模式和 Fisher 信息定义列出来；“两点弱非相干源”的结论只在相应条件下使用。

\section{天体辐射机制和源模型}
\label{appsec:g-sources}

恒星、maser、自然激光、Crab 光子统计和紧致天体章节需要同时读经典辐射机制和现代观测。量子光学语言无法替代普通天体物理模型；它的作用是标出哪些光子统计量还没有被平均掉。读这一组文献时，按源类型记录三个量：亮度或光子率、角尺度或时间尺度、以及最可能破坏理想统计的混合成分。对受激辐射和自然激光候选，尤其要区分窄线、高亮温、泵浦涨落和稳定相干态之间的差别 \citep{1982RvMP...54.1225E,1992ARA&A..30...75E,2004A&A...428..497J,2005MNRAS.364..731J,2005NewA...10..361J,2008AIPC..984..216D}。

\section{传播、透镜、偏振和宇宙学}
\label{appsec:g-propagation}

传播章节的文献应按观测量来读：DM 和色散延迟、RM 和偏振角、散射脉冲展宽、尘埃消光、引力透镜时间延迟、波动光学透镜、CMB 偏振和宇宙双折射。每个主题都要问同一个问题：它改变的是到达时间、频率、偏振、相位，还是强度相关？如果这个问题答不清，后文很容易把传播效应误写成源的量子统计。

\section{科学案例和远期网络}
\label{appsec:g-cases}

第一代科学案例可以从亮星角直径、快速自转星、Be 星盘、双星、Crab 光子统计、自然激光和近邻瞬变开始。Type Ia 超新星距离方案展示了强度干涉怎样进入宇宙学距离问题，但同时也暴露了基线、光子率、触发时间和模型系统误差的压力 \citep{2025PhRvD.111h3047K}。量子网络望远镜则应从 pathfinder 读起：two-photon astrometry、纠缠辅助干涉、存储和频率转换的资源约束分别见 \citet{2023PhRvL.130p0801M,2025PhRvA.111d3701M,2025PhRvR...7d3278Z,2026PhRvA.113a2608P,2026Natur.651..326S}。这组文献适合放在课程最后，因为事件表、相干函数、仪器误差和误差预算先熟悉以后，网络方案的资源压力才会清楚。

\section{读文献时的记录模板}
\label{appsec:g-reading-template}

\begin{longtable}{p{0.22\linewidth}p{0.68\linewidth}}
\toprule
记录项 & 写法 \\
\midrule
观测量 & 写出论文实际测的量，例如 \(g^{(2)}(\tau)\)、\(|V|^2\)、角直径、偏振角、延迟或事件率。 \\
数据产品 & 记录论文是否发布事件表、相关直方图、校准表、后验样本或代码。 \\
主要数量级 & 光子率、带宽、时间分辨率、基线、目标星等、背景和积分时间。 \\
核心假设 & 热光、高斯 PSF、均匀圆盘、非相干源、稳定响应、弱信号等。 \\
系统误差 & 背景、死时间、afterpulsing、串扰、校准星、模型误差和选择效应。 \\
可复用图像 & 论文中哪些图可以启发教学图，但不能照搬。 \\
\bottomrule
\end{longtable}
